\chapter{Практический раздел}

\section{Структура приложения}

Проект реализован на языке Python и состоит из набора модулей, разделенных по
ответственности. В модуле модели предметной области описаны сущности
\texttt{Tour}, \texttt{Client} и \texttt{Booking}. Отдельно выделены сервисы
каталога, клиентов и бронирований. Подсистема маркетинга реализована в модуле
\texttt{LegacyMarketingHub}, а подсистема уведомлений --- через издателя и
наблюдателей.

Пользовательское взаимодействие организовано через консольный интерфейс.
Демонстрационный сценарий позволяет просматривать каталог туров, регистрировать
клиентов, создавать бронирования и получать маркетинговый отчет. Такой способ
запуска достаточен для учебной работы и позволяет проверить ключевые сценарии
без дополнительной инфраструктуры.

\section{Реализация ключевых сценариев}

\subsection{Создание бронирования}

При создании бронирования выполняется следующая последовательность действий:

\begin{enumerate}
    \item проверяется существование клиента и тура;
    \item рассчитывается итоговая стоимость с учетом скидок;
    \item резервируются места в туре;
    \item создается объект бронирования;
    \item публикуются уведомления о создании брони и наличии задолженности.
\end{enumerate}

Если во время выполнения операции возникает ошибка, Unit of Work откатывает
резерв мест и удаляет промежуточно созданные данные. За счет этого система
сохраняет согласованность даже при частичном отказе.

\subsection{Управление жизненным циклом брони}

Сервис бронирования поддерживает перевод заявки в работу, отметку о готовности,
фиксацию оплаты, выдачу и отмену. Ограничения на переходы определяются
состояниями объекта бронирования. Например, выдача разрешается только после
полной оплаты, а отмена после статуса <<Готова>> запрещена. Такой подход
делает бизнес-правила явными и уменьшает вероятность некорректных изменений
статуса.

\subsection{Маркетинговая аналитика}

Маркетинговый модуль анализирует историю бронирований клиента, суммарные
расходы, задолженность и предпочтительные направления. На этой основе клиент
относится к одному из сегментов: \texttt{welcome}, \texttt{reactivation},
\texttt{vip}, \texttt{winback}, \texttt{weekend} или \texttt{cross\_sell}.
Далее формируются канал коммуникации, рекомендованный тур и промокод. Такой
сценарий показывает, как существующую бизнес-логику можно изолировать в
самостоятельном модуле и использовать без прямой зависимости от интерфейса.

\section{Тестирование}

Для проверки корректности реализованы модульные тесты. Они покрывают следующие
сценарии:

\begin{itemize}
    \item допустимые переходы состояний и публикацию уведомлений;
    \item запрет отмены после перевода брони в состояние готовности;
    \item откат зарезервированных мест при ошибке сохранения;
    \item однократную отправку уведомления об истечении срока действия брони;
    \item формирование изолированного маркетингового отчета для разных типов клиентов.
\end{itemize}

Такой набор тестов проверяет не только отдельные функции, но и взаимодействие
между сервисами, репозиториями и паттернами.

\section{Результат работы}

В ходе лабораторной работы была получена учебная система туристического
агентства, в которой выделены основные сущности предметной области, реализованы
гибкие правила ценообразования, формализован жизненный цикл бронирования и
добавлены уведомления с маркетинговой аналитикой. Архитектура проекта
подготовлена к дальнейшему расширению и может быть адаптирована для подключения
реальной базы данных или сетевого интерфейса~\cite{python312}.
